\section{Introduction}
\label{s:intro}

Panels of reported cases by county or region are standard for evaluating non-pharmaceutical
interventions such as school closures, mobility restrictions and community mandates. In practice,
effects are commonly estimated by event studies, two-way fixed-effects regressions of case growth
or log cases, regional comparisons and synthetic controls
\citep{hsiang2020anticontagion,chernozhukov2021masks,mitze2020masks}, importing the
difference-in-differences logic of policy evaluation
\citep{callaway2021did,roth2023trending,haber2021impact}. Epidemic data strain that logic on
three fronts simultaneously, because transmission generates the dependence that outcome-scale
designs treat as error. Weekly cases arise from earlier cases through the generation interval
\citep{wallinga2004different,fraser2007estimating,wallinga2007generation}, so parallel trends in
counts or log counts fail even under no policy \citep{gauthier2021twfe,feng2024parallel}. Policy
adoption often responds to rising incidence, invalidating strict exogeneity and conventional unit
fixed effects \citep{imai2019fixed,zhang2025timevarying}. Finally, units are spatially coupled:
infections cross borders and neighbours' interventions alter contact patterns
\citep{hudgens2008interference,halloran2016dependent,aronow2017exposure}, biasing estimates when
spillovers are omitted \citep{enriquez2025spillovers}.

Existing approaches each resolve one piece of this structure. \citet{feng2024parallel} place
parallel trends on cohort-based log effective reproduction numbers, omitting spillovers and relying
on first-step reproduction-number estimates subject to small-count Jensen bias and
generation-interval distortions \citep{gostic2020practical}. Semi-mechanistic renewal models in
Bayesian \citep{flaxman2020npi,brauner2021npi,bhatt2023semimech} or frequentist form
\citep{bong2024frequentist,bong2024augmented}, SIR-motivated unconfoundedness
\citep{callaway2023pandemic} and structural growth models \citep{chernozhukov2021masks} condition on
the observed past rather than eliminating unrestricted unit effects. Spatio-temporal count models
\citep{held2005hhh,meyer2014powerlaw,alaimo2024poissonar} capture contagion coupling without
incidence-driven policy feedback, whereas spatial and network difference-in-differences designs
\citep{kosfeld2021spatial,butts2021spatial,forastiere2021network,xu2023did,sun2025didnetwork}
model spillovers on the outcome scale without renewal dynamics. Quasi-differenced sequential
moments remove multiplicative unit effects under predetermined regressors
\citep{chamberlain1992sequential,wooldridge1997multiplicative,windmeijer2000moment,blundell2002count}
but have not been integrated with renewal transmission, spatial contagion and formal validity
diagnostics.

This paper develops quasi-differenced generalized method of moments (QD-GMM) on the transmission
scale together with explicit identification and failure-mode diagnostics. Rather than assuming
weekly surveillance panels can always recover structural causal effects, we characterize precisely
when transmission parameters are point-identified and when diagnostics downgrade estimates to
pseudo-true values of a rejected working model. Parallel trends are placed on the log reproduction
number of a renewal equation whose infection pressure imports cases from neighbouring units; a
forward quasi-difference eliminates unrestricted unit effects; instruments built solely from the
observed past permit incidence-driven feedback; and policy spillovers enter through the share of
treated neighbours. Three methodological results follow. First, under cross-cluster conditions
excluding unmodelled same-week cross-cluster transmission, cluster scores form a martingale
difference array across weeks despite lagged inter-cluster contagion, so the ordinary
cluster-robust sandwich applies without spatial HAC adjustment and week effects profile out in
closed form. Second, across pre-specified simulations with feedback adoption, spillover,
overdispersion and small counts ($R=8000$), QD-GMM recovers direct effects with bias between
$-0.002$ and $0.001$ and 95\% coverage of 94.5\% to 95.4\% where outcome-scale, Poisson
fixed-effects and two-step log-$R$ comparators fail, while boundary ($\phi_0=0$) and same-week
($w_0=0.3$) experiments demarcate the estimator's domain of validity. Third, built-in
diagnostics---profiled contagion curves that reveal multiple roots, over-identification tests
($J$ or rank-safe score-form $J_{LM}$) calibrated against finite-sample null distributions,
Stock--Wright identification-robust sets, total-policy-scale event-time profiles and
ascertainment-confounding checks---detect specification and surveillance failures directly from
observed data.

We stress-test this framework on five years of German county-level weekly varicella
(\emph{Windpocken}) notifications (400 counties, 16 states, 2015--2019; 115{,}129 cases) paired
with federal-state school holiday calendars
\citep{cauchemez2008school,jackson2014varicella,adda2016economic,vonbismarckosten2022schools,isphording2021schools}.
Varicella's 14-day mean serial interval \citep{vink2014serial} keeps within-week transmission
negligible ($w_0=2.6\times 10^{-6}$), while sparse county counts (58.8\% zeros) test Theorem~3's
predictions for standard comparators. On the empirical data, two-step EpiEstim log-$R$ regression
flips sign to $+0.064$ before reverting to $-0.102$ at NUTS-2 aggregation (2.9\% zeros), and
$\log(Y+1)$ TWFE attenuates to $-0.102$. Under the baseline QD-GMM working model, the policy-level
contrast is negative ($\hat\tau+\hat\kappa=-0.239$, 95\% $t_{15}$ CI $-0.364$ to $-0.115$;
restricted score wild bootstrap $-0.419$ to $-0.139$) with an interior coupling root
$\hat\phi=0.326$. Yet the framework's diagnostics prevent taking this working-model contrast as an
unconfounded 21.3\% causal reduction: $J_{LM,\mathrm{oi}}$ rejects lagged working-model moments
($p_{\mathrm{cal}}=0.004$), total-scale event-time estimates show an immediate $-0.208$ onset-week
drop ($p=0.049$) inconsistent with notification weeks and a 14-day incubation lag, short non-summer
holidays display the classic reporting-suppression and lag-2 rebound echo, the two largest
southern states (Baden-W\"urttemberg and Bavaria) never rotate their start timing, and dropping
both states or absorbing four-region $\times$ week effects yields intervals covering zero.
Sections~\ref{s:methods}--\ref{s:discussion} present the theory, simulations, German stress test
and scope; Web Appendices~A--F provide proofs, adapted comparator results, numerical checks,
supplementary simulations, data provenance, few-cluster calibration and full diagnostics.
